%% ma: L11-B04-0013
%% lop: 11
%% chuong: 1
%% bai: 4
%% dang: 11.4.1
%% loai: DS
%% muc: [TH, NB, TH, VD]
%% dap_an: "ĐSSĐ"
%% nguon: "de-cuong-giua-ky-1-toan-11-nam-2026-2027-truong-thpt-hoang-van-thu-ha-noi-2-8.pdf (M08, Đề cương giữa kì 1 Toán 11 THPT Hoàng Văn Thụ, Hà Nội), Đề minh họa, Phần II Câu 2"
%% kien_thuc: "Bài 3 (hàm số lượng giác), Bài 4 (phương trình lượng giác cơ bản)"
%% trang_thai: chinh-thuc
\begin{ex}
Hàng ngày mực nước của một con kênh lên xuống theo thuỷ triều. Độ sâu $h$ (mét) của mực nước trong kênh tính theo thời gian $t$ (giờ) được cho bởi công thức $h(t)=3\cos\left(\dfrac{\pi t}{12}+\dfrac{\pi}{3}\right)+12$ với $0\le t\le24$.
\choiceTF{\True Độ sâu của mực nước trong kênh lớn nhất là $15$ m}{Mực nước trong kênh tại thời điểm $t=6$ giờ là $10$ m}{Có đúng hai thời điểm trong ngày mực nước của kênh có độ sâu $9$ m}{\True Thời điểm $8$ giờ tối là lúc độ sâu của mực nước trong kênh đạt giá trị lớn nhất}
\loigiai{a) Đúng: $h(t)\le3+12=15$, đạt tại $t=20$.

b) Sai: $h(6)=3\cos\dfrac{5\pi}6+12=12-\dfrac{3\sqrt3}2\approx9{,}4$ (m).

c) Sai: $h=9\Leftrightarrow\cos\left(\dfrac{\pi t}{12}+\dfrac\pi3\right)=-1\Leftrightarrow t=8+24k$; trên $[0;24]$ chỉ có $t=8$.

d) Đúng: $8$ giờ tối ứng với $t=20$, $h(20)=3\cos2\pi+12=15$ là giá trị lớn nhất.}
\end{ex}
